\section{Analyse der angewendeten Forschungsmethoden}

\subsection{Quantitative Methoden der Studie}

\subsubsection{Datenerhebung und Fragebogen}

Die quantitative Datenerhebung der Studie ist eine teil-standardisierte Online-Befragung. Beworben wurde die Online-Befragung von zahlreichen Organisationen, Schulen, Kitas, Fachkräften und sonstigen Multiplikatoren. Teilnehmende, die nicht den Kriterien der Zielgruppe entsprechen, wurden durch eine entsprechende Führung im Fragebogen aussortiert.

Die Konstruktions des Fragebogens geschah auf Basis der Literatur- und Datenanalyse und der Ergebnisse der Fokusgruppe aus Haupt- und Ehrenamtlichen, welche im Bereich der Prävention tätig sind. Vor Durchführung der Online-Befragung wurde der Fragebogen von Mitarbeitenden der Stadtverwaltung und Mitgliedern des Runden Tischs in einem Pretest getestet und entsprechend angepasst. Der Fragebogen enthält offene, halb-offene und geschlossene Fragen.

\subsubsection{Stichprobe}

Die Grundgesamtheit bilden alle Familien in Tübingen mit Kindern unter 18 Jahren, die erstens geringes Einkommen oder zweitens Sozial- oder Transferleistungen beziehen oder drittens knapp über der Bezugsgrenze liegen. Insgesamt haben im Zeitraum von fünf Monaten 355 Familien teilgenommen. Das Ausfüllen des Fragebogens wurde dabei von einem Elternteil oder einer assistierenden Person vorgenommen.

Die Stichprobe bildet eine Quotenwahl aus Gruppen, welche nach für die Autoren armutsrelevanten Kriterien eingeteilt wurden. Darunter fallen etwa die Familiensituation, der Bezug verschiedener Sozialleistungen, die Erwerbssituation und der Wohnort. Nicht alle Quoten wurden erreicht, weshalb die Autoren erklären, dass keine Rückschlüsse auf die Gesamtbevölkerung möglich sind. Dass jedoch alle Gruppen mit mindestens 30 Befragten vertreten sind, bietet die Möglichkeit relevante Anhaltspunkte zu erkennen und Gruppen untereinander zu vergleichen.

\subsubsection{Operationalisierung}

TODO?

\subsubsection{Datenauswertung}

Vor Auswertung wurden die Daten aufbereitet und bereinigt. Geschlossene Fragen wurden statistisch analysiert, wobei unter anderem deskriptive und bivariante Analysen durchgeführt wurden. Die Antworttexte der offenen und halb-offenen Fragen wurden durch eine inhaltlich strukturierende qualitative Inhaltsanalyse angelehnt an Mayring (2016 TODO) analysiert. Die Kategorien wurden dabei induktiv aus den Daten entwickelt und ihre Nennungen quantifiziert.

\subsection{Qualitative Methoden der Studie}

\subsubsection{Fokusgruppen}

Insgesamt wurden drei Fokusgruppen aus Familien bzw. Elternteilen mit 22 Teilnehmenden und drei Fokusgruppen aus haupt- und ehrenamtlichen Unterstützern mit 22 Teilnehmenden. Erwähnt wird weiterhin ein studentisches Teilprojekt aus zwei Fokusgruppen und einem Einzelinterview mit Jugendlichen am Campus Reutlingen.

Ausgewertet wurden die Fokusgruppendiskussionen mittels einer qualitativen Inhaltsanalyse angelehnt an Kuckartz und Rädiker (2022, TODO). Auf die Transkription nach Dresing und Pehl (2018, TODO) folgte die inhaltsanalytische Auswertung der einzelnen Fokusgruppen und eine anschließendes Verknüpfen der Erkenntnisse aus den jeweiligen Fokusgruppen.

\subsubsection{Literatur- und Datenanalyse und Workshop}

Weitere Methoden neben der quantitativen Online-Befragung und den qualitativen Fokusgruppen waren die Literatur- und Datenanalyse sowie der Workshop des Runden Tischs, welche jedoch eine untergeordnete Rolle spielen. Die Literatur- und Datenanalyse bestand aus einer Sichtung von Konzeptunterlagen und Informationsmaterial der Angebote des Tübinger Präventionskonzepts. Aus der systematischen Angebotsanalyse wurden fünf Handlungsfelder herausgearbeitet, welche als Informationsgrundlage für die Entwicklung des Fragebogens und der Fokusgruppen-Leitfäden eingesetzt wurden. Der Workshop des Runden Tischs gegen Ende der Studie beinhaltete eine Präsentation und Diskussion der bisherigen Evaluationsergebnisse. Ausgehend davon wurden weitere Projektideen für die Handlungsfelder überlegt und hinsichtlich ihrer Durchführbarkeit geprüft.